%%%%%%%% mlsys 2025 EXAMPLE LATEX SUBMISSION FILE %%%%%%%%%%%%%%%%%

\documentclass{article}

% Recommended, but optional, packages for figures and better typesetting:
\usepackage{microtype}
\usepackage{graphicx}
\usepackage{subfigure}
\usepackage{booktabs} % for professional tables
\usepackage{amsmath}

% hyperref makes hyperlinks in the resulting PDF.
% If your build breaks (sometimes temporarily if a hyperlink spans a page)
% please comment out the following usepackage line and replace
% \usepackage{mlsys2025} with \usepackage[nohyperref]{mlsys2025} above.
\usepackage{hyperref}

% Attempt to make hyperref and algorithmic work together better:
\newcommand{\theHalgorithm}{\arabic{algorithm}}

% Use the following line for the initial blind version submitted for review:
% \usepackage{mlsys2025}

% If accepted, instead use the following line for the camera-ready submission:
\usepackage[accepted]{mlsys2025}

% The \mlsystitle you define below is probably too long as a header.
% Therefore, a short form for the running title is supplied here:
\mlsystitlerunning{Optimizing DSA GPU Kernel for LLM Inference Using Triton}

\begin{document}

\twocolumn[
\mlsystitle{Optimizing DeepSeek Sparse Attention GPU Kernel for LLM Inference Workloads Using Triton}

% It is OKAY to include author information, even for blind
% submissions: the style file will automatically remove it for you
% unless you've provided the [accepted] option to the mlsys2025
% package.

% List of affiliations: The first argument should be a (short)
% identifier you will use later to specify author affiliations
% Academic affiliations should list Department, University, City, Region, Country
% Industry affiliations should list Company, City, Region, Country

% You can specify symbols, otherwise they are numbered in order.
% Ideally, you should not use this facility. Affiliations will be numbered
% in order of appearance and this is the preferred way.
\mlsyssetsymbol{equal}{*}

\begin{mlsysauthorlist}
\mlsysauthor{Mengze Tang}{to}
\end{mlsysauthorlist}

\mlsysaffiliation{to}{Ph.D. Student, Department of Computer Sciences, University of Wisconsin-Madison, Madison, WI, United States}

\mlsyscorrespondingauthor{Mengze Tang}{mengze@cs.wisc.edu}
% \mlsyscorrespondingauthor{Eee Pppp}{ep@eden.co.uk}

% You may provide any keywords that you
% find helpful for describing your paper; these are used to populate
% the "keywords" metadata in the PDF but will not be shown in the document
\mlsyskeywords{Machine Learning, MLSys}

\vskip 0.3in

\begin{abstract}
DeepSeek Sparse Attention (DSA) reduces the main attention computation in long-context
LLM inference by selecting a fixed number of key-value entries per query token before
running attention. This report studies the kernel-level part of that computation: given
the selected token indices, compute the DSA attention output and log-sum-exp values for
the DeepSeek-V3.2-Exp-style MLA/MQA layout. Starting from a PyTorch reference
implementation, I developed three Triton proposals that progressively moved the
computation into fused GPU kernels, changed the tiling strategy, and specialized the
positional-key path. On the B200 benchmark suite, the best version is correct on all 23
workloads and reaches a reported mean speedup of $4.74\times$ with average latency
reduced from about $1.62$ ms to $0.40$ ms.
\end{abstract}
]

% this must go after the closing bracket ] following \twocolumn[ ...

% This command actually creates the footnote in the first column
% listing the affiliations and the copyright notice.
% The command takes one argument, which is text to display at the start of the footnote.
% The \mlsysEqualContribution command is standard text for equal contribution.
% Remove it (just {}) if you do not need this facility.

\printAffiliationsAndNotice{}  % leave blank if no need to mention equal contribution
% \printAffiliationsAndNotice{\mlsysEqualContribution} % otherwise use the standard text.

\section{Problem Statement / Introduction}
\label{submission}

The initial project proposal targeted optimized attention kernels for LLM inference,
with sparse attention as the main case study. The motivation was that attention becomes
one of the dominant costs in long-context inference: dense causal attention grows
quadratically with sequence length, and even optimized dense kernels must move and
combine a large amount of KV-cache data. The proposal therefore planned to study
DeepSeek Sparse Attention (DSA), implement relevant Triton kernels, and benchmark
them against a reference implementation.

The final project scope narrowed in a useful way. Instead of implementing the full DSA
pipeline, including the lightning indexer and top-$k$ selector, this work focuses on the
core sparse-attention kernel after the selected token indices are already available. That
is the part represented by \texttt{baselines/attention.py} and the task
\texttt{dsa\_sparse\_attention}, with 16 query heads, 512 compressed-cache dimensions,
64 positional-key dimensions, 2048 selected tokens, and 64-token pages. The upstream
indexer determines which tokens should be attended to; the kernel studied here must
efficiently gather those tokens from the paged cache, compute attention logits for all
query heads, produce a numerically stable softmax, and return both the BF16 output and
base-2 log-sum-exp (LSE).

\subsection{DSA Background}

DeepSeek-V3.2-Exp introduces DSA as a sparse-attention modification to
DeepSeek-V3.1-Terminus \cite{deepseekv32}. The DSA design has two conceptually
separate components. First, a lightweight indexer assigns an index score $I_{t,s}$
between a query token $t$ and each preceding token $s$. Second, a fine-grained selector
keeps only the key-value entries with top-$k$ index scores, and the main attention
computation runs over this selected set. In the sparse training stage described in the
DeepSeek-V3.2-Exp report, $k=2048$, matching the benchmark used in this project.

This architecture changes the main attention complexity from $O(L^2)$ to $O(Lk)$ for
sequence length $L$, while leaving a much smaller indexer computation upstream. The
paper also instantiates DSA under Multi-head Latent Attention (MLA) \cite{deepseekv2}
in an MQA-like mode \cite{shazeer2019}, where one latent key-value entry is shared
across multiple query heads. That detail is directly reflected in the benchmark layout:
each query token has 16 query heads, but the paged cache stores one 512-dimensional
compressed key/value vector and one 64-dimensional positional key vector per cached
token. This is also consistent with hardware-aware sparse-attention work that emphasizes
aligning sparse patterns with GPU execution constraints \cite{yuan2025nsa}.

\subsection{Kernel Interface}

For each query token $t$ and head $h$, the reference implementation computes over
the selected token set $S_t$ encoded by \texttt{sparse\_indices[t]}. A selected global
token id $s$ maps to page $p=\lfloor s/64\rfloor$ and page offset $r=s\bmod 64$.
The benchmark constants are:
\[
\begin{aligned}
H &= 16,\quad d_c = 512,\quad d_r = 64, \\
k &= 2048,\quad \text{page size} = 64.
\end{aligned}
\]
The logit for one selected token is
\[
\ell_{t,h,s} =
\alpha\left(
q^c_{t,h}\cdot c_{p,r} + q^r_{t,h}\cdot k^r_{p,r}
\right),
\]
where $\alpha$ is the softmax scale, $c_{p,r}$ is the compressed latent cache entry
used as both key content and value, and $k^r_{p,r}$ is the positional key. The output is
\[
o_{t,h}=\sum_{s\in S_t}
\frac{\exp(\ell_{t,h,s})}{\sum_{u\in S_t}\exp(\ell_{t,h,u})}c_{p(s),r(s)}.
\]
The LSE value returned by the benchmark is
\[
\mathrm{LSE}_{t,h}=\log_2\sum_{s\in S_t}\exp(\ell_{t,h,s}).
\]

\section{Solution Description}

The baseline in \texttt{baselines/attention.py} is a PyTorch implementation with a
token-level Python loop. For each token, it filters invalid \texttt{-1} entries, gathers
the selected rows from the flattened paged caches, computes the two logit terms with
matrix multiplications, runs \texttt{torch.logsumexp} and \texttt{torch.softmax}, and
then writes the BF16 output. From a performance perspective, this materializes
intermediate tensors such as the gathered keys, logits, and probabilities. The repeated
gather and softmax work across 16 heads and 2048 selected tokens is the main target for
GPU fusion.

All three proposals keep the same high-level two-pass algorithm:

\begin{enumerate}
\item A stage-1 Triton kernel scans the sparse token set for each $(t,h)$ pair and
computes a numerically stable softmax summary: the maximum logit $m$ and the
normalizer $\sum_s\exp(\ell_s-m)$. It also writes the required base-2 LSE.
\item A stage-2 Triton kernel rescans the same sparse token set, recomputes logits,
normalizes them using the saved $m$ and normalizer, and accumulates the weighted
512-dimensional value vector.
\end{enumerate}

The two-pass structure intentionally recomputes logits in stage 2. This avoids storing a
large probability or logit tensor of shape
\texttt{num\_tokens} $\times 16 \times 2048$, which would be expensive in both
memory footprint and bandwidth. The cost is an additional read of the selected cache
entries, but the benchmark shape is small enough per query/head that reducing global
intermediate storage is the better design point.

\subsection{Proposal 1: Functional Two-Stage Triton Kernel}

The first proposal moved the sparse attention work into Triton. Stage 1 launches a grid
of \texttt{(num\_tokens, 16)} programs, one program per token and query head. It scans
the 2048 selected indices in blocks of 32, gathers the corresponding compressed cache
and positional cache vectors through page/offset arithmetic, and updates the online
softmax statistics. Stage 2 launches
\texttt{(num\_tokens, 16, 8)} programs, where the third dimension splits the
512-dimensional output into eight 64-element value tiles.

This proposal establishes the important optimization boundary: the Python loop,
intermediate gather tensors, and PyTorch softmax tensor are removed from the measured
kernel path. The online LSE update uses base-2 exponentials internally and converts
natural-logits through $\log_2(e)$, which matches the benchmark's base-2 LSE
contract. The first proposal used conservative tile sizes
\texttt{BLOCK\_K=32}, \texttt{BLOCK\_DCKV=32}, and \texttt{BLOCK\_DKPE=32}. Those
settings made compilation and correctness easier, but they also produced many loop
iterations: 64 sparse-index tiles per pass and 16 content-dimension tiles inside each
logit computation.

\subsection{Proposal 2: Larger K and Feature Tiles}

The second proposal kept the same two-stage decomposition but changed the tile shape.
It increased \texttt{BLOCK\_K} from 32 to 64 and increased the content and positional
feature tiles to 64. This halves the number of sparse-index iterations per pass and also
halves the number of content-vector tiles needed to compute each logit block. Because
the positional key dimension is exactly 64, \texttt{BLOCK\_DKPE=64} also reduces the
positional path to one tile.

The kernel body was also simplified around this tile shape: the unused \texttt{num\_pages}
argument was removed from the JIT kernels, the logarithm conversion factor was stored
as a local scalar, invalid entries were masked before contributing to the softmax sum,
and the wrapper normalized tensor contiguity and scalar \texttt{sm\_scale} handling
before launch. The performance-relevant change is the tile enlargement: it reduces loop
control, pointer construction, and repeated query loads while exposing a larger vector of
selected tokens to each Triton program.

\subsection{Proposal 3: KPE Specialization and Larger Sparse Tiles}

The third proposal specialized the fixed 64-dimensional positional-key computation and
increased \texttt{BLOCK\_K} from 64 to 128. Since $d_r=64$ is a benchmark constant, the
generic loop over positional-key tiles was replaced by a single vectorized load and dot
product in both stage 1 and stage 2. This removes loop overhead and gives the compiler a
more direct representation of the small positional term.

The larger \texttt{BLOCK\_K=128} tile reduces the number of sparse-index blocks from 32
to 16 per pass. It also amortizes query-vector loads and softmax-update overhead over
more selected tokens. The improvement is smaller than the Proposal 1 to Proposal 2 jump
because the larger sparse tile increases the live working set inside each program; at
some point, fewer loop iterations must be traded against register pressure and memory
traffic from the gathered cache rows.

\section{Overview of Results / Evaluation}

The experiments were evaluated with the Modal backend described in
\texttt{modal\_eval.py}. The configuration uses Triton on an NVIDIA B200 GPU, with
23 benchmark workloads for the DSA sparse-attention task. Each result below compiled,
passed numerical correctness, and was compared against the evaluator's reference
latency for the same task. The reported speedup is the benchmark's mean of per-workload
speedup factors, so it is not exactly the ratio of the average reference latency to the
average proposal latency.

\begin{table}[h]
\centering
\caption{DSA sparse-attention benchmark results on B200.}
\label{tab:results}
\resizebox{\columnwidth}{!}{%
\begin{tabular}{lrrrr}
\toprule
Version & Latency (ms) & Speedup & Abs. err. & Rel. err. \\
\midrule
Proposal 1 & 0.803 & $2.19\times$ & 0.01023 & 0.0361 \\
Proposal 2 & 0.483 & $3.80\times$ & 0.00435 & 0.0442 \\
Proposal 3 & 0.401 & $4.74\times$ & 0.00306 & 0.0341 \\
\bottomrule
\end{tabular}%
}
\end{table}

Proposal 1 already gives a substantial improvement because it changes the execution
model from many PyTorch operations around gathered tensors into two GPU kernels with
explicit tiling. Its reported mean speedup is $2.19\times$, with average latency
$0.803$ ms compared with an average reference latency of about $1.62$ ms. The remaining
cost comes from the small tile sizes: stage 1 and stage 2 each perform many short loops
over both selected-token blocks and the 512-dimensional compressed cache.

Proposal 2 is the largest single step. Average latency falls from $0.803$ ms to
$0.483$ ms, a $39.8\%$ reduction, and reported mean speedup rises from $2.19\times$ to
$3.80\times$. The reason is consistent with the kernel structure: \texttt{BLOCK\_K=64}
reduces sparse-index loop iterations by $2\times$, and \texttt{BLOCK\_DCKV=64} reduces
the number of content-vector chunks by $2\times$. Fewer loop iterations mean fewer
pointer calculations, fewer masked loads of query slices, and better amortization of
the online softmax update.

Proposal 3 improves average latency further to $0.401$ ms and reaches a reported
$4.74\times$ mean speedup. Relative to Proposal 2, this is a $17.1\%$ latency reduction.
The fixed-size positional term is a good specialization target because every selected
token must include the 64-dimensional positional-key dot product, but that dimension
never changes across workloads. Increasing the sparse tile to 128 also reduces scan
overhead. The more modest gain suggests that the kernel is approaching a balance point:
larger tiles reduce control overhead but make each program heavier, especially because
the selected cache rows are still gathered indirectly through sparse page indices.

Overall, the best kernel reduces average latency by about $75.2\%$ relative to the
average reference latency in the result file, while preserving correctness on all 23
workloads. The result also illustrates an important sparse-attention lesson: once the
algorithm has reduced the token set to a fixed top-$k$, the kernel bottleneck becomes
efficient indirect memory access and careful tiling, not the asymptotic attention
formula alone.

\section{Deliverables}

\section{Conclusions and Future Work}

This project started with a broad plan to study sparse attention, top-$k$ selection, and
Triton implementation techniques. The final work concentrated on the DSA core-attention
kernel, which was the right level of scope for making concrete GPU-performance progress.
The main result is an iterative Triton implementation that moves from a correct
two-stage sparse-attention kernel to a more specialized tile shape, improving reported
mean speedup from $2.19\times$ to $4.74\times$.

The most useful optimization lesson was that the first correct GPU mapping is not enough
for this workload. The sparse token set is fixed at 2048 entries, and each entry requires
a 512-dimensional content dot product plus a 64-dimensional positional dot product.
Small tiles spend too much time in loop overhead and repeated address formation. Larger
tiles improve utilization, but only up to the point where register pressure and indirect
cache loads begin to dominate. Specializing known constants such as the 64-dimensional
positional key is therefore valuable because it removes overhead without changing the
algorithm.

Future work should profile the best kernel with hardware counters to separate memory
bandwidth, L2 behavior, register pressure, and occupancy effects. The current two-pass
design avoids storing logits, but it recomputes the content and positional dot products;
more advanced variants could split the selected-token dimension across programs and use
a reduction, or cache partial statistics when the extra storage is justified. Another
direction is to optimize the upstream lightning indexer in \texttt{baselines/indexer.py}
and study the full DSA pipeline, including FP8 dequantization and top-$k$ selection.
Finally, this work directly uses ME759 topics: GPU thread/block decomposition, memory
hierarchy awareness, tiling, performance benchmarking, and parallel algorithm design for
modern ML inference.

% Acknowledgements should only appear in the accepted version.
% \section*{Acknowledgements}


\begin{thebibliography}{9}

\bibitem[DeepSeek-AI(2025)]{deepseekv32}
DeepSeek-AI.
\newblock DeepSeek-V3.2-Exp: Boosting Long-Context Efficiency with DeepSeek Sparse Attention.
\newblock Technical report, 2025.

\bibitem[DeepSeek-AI(2024)]{deepseekv2}
DeepSeek-AI.
\newblock DeepSeek-V2: A Strong, Economical, and Efficient Mixture-of-Experts Language Model.
\newblock arXiv:2405.04434, 2024.

\bibitem[Shazeer(2019)]{shazeer2019}
Noam Shazeer.
\newblock Fast Transformer Decoding: One Write-Head is All You Need.
\newblock arXiv:1911.02150, 2019.

\bibitem[Yuan et~al.(2025)]{yuan2025nsa}
Jingyang Yuan, Huazuo Gao, Damai Dai, and others.
\newblock Native Sparse Attention: Hardware-Aligned and Natively Trainable Sparse Attention.
\newblock Proceedings of ACL 2025.

\end{thebibliography}


\end{document}


% This document was modified from the file originally made available by
% Pat Langley and Andrea Danyluk for ICML-2K. This version was created
% by Iain Murray in 2018. It was modified from a version from Dan Roy in
% 2017, which was based on a version from Lise Getoor and Tobias
% Scheffer, which was slightly modified from the 2010 version by
% Thorsten Joachims & Johannes Fuernkranz, slightly modified from the
% 2009 version by Kiri Wagstaff and Sam Roweis's 2008 version, which is
% slightly modified from Prasad Tadepalli's 2007 version which is a
% lightly changed version of the previous year's version by Andrew
% Moore, which was in turn edited from those of Kristian Kersting and
% Codrina Lauth. Alex Smola contributed to the algorithmic style files.
